\documentclass[11pt]{article}
\usepackage[margin=1in]{geometry}
\usepackage{booktabs}
\usepackage{hyperref}
\usepackage{xcolor}
\usepackage{longtable}
\usepackage{array}
\title{Replication Report: Fang et al.\ (2018)\\
\large \emph{E. coli B2 strains prevalent in IBD patients have distinct metabolic capabilities that enable colonization of intestinal mucosa}}
\author{Ollie (OpenClaw AI) --- BVBRC Replication Project, Wave 4, target \#17}
\date{2026-06-27}
\begin{document}
\maketitle

\section*{Bibliographic anchor}
Fang X, Monk JM, Mih N, Du B, Sastry AV, Kavvas E, Seif Y, Smarr L, Palsson BO.
\emph{BMC Systems Biology} 12:66 (2018).
DOI: \href{https://doi.org/10.1186/s12918-018-0587-5}{10.1186/s12918-018-0587-5};
PMC5996543; PMID 29890970. Open access (CC BY 4.0 / BMC).

\section*{Verdict}
\textbf{PARTIAL REPLICATION (strong).} Both the paper's central FBA
differentiation claim (Table 1) and its mechanistic loss-of-function
claim (C5: B2 lacks the \emph{frl} operon) are independently reproduced
on actual public data: FBA on the K-12 reference GEM and \texttt{tblastn}
on the three canonical B2 reference genomes (LF82, UTI89, NRG857c).
Phylogroup assignments were re-derived via in-silico Clermont quadruplex
PCR. The 4-strain genomic panel still falls short of the paper's
110-strain pan-genome rerun, so the verdict is PARTIAL rather than full
REPLICATED.

\section{Paper claims tested}
\begin{longtable}{clp{5.5cm}c}
\toprule
\# & Claim & Type & Tested? \\
\midrule
C1 & Enough public \emph{E. coli} genomes exist to reconstruct the 110-strain pan-genome & Data availability & \checkmark \\
C2 & LF82, UTI89, NRG857c publicly available; genome stats match & Data + stats & \checkmark \\
C3 & B2 has metabolic genes differentially distributed vs.\ other phylogroups (loss-of-function) & Genomic & Partial (4 strains) \\
C4a & Non-B2 GEMs predict growth on Amadori products (fructoselysine, psicoselysine) & FBA & \checkmark \\
C4b & Non-B2 grows on Table-1 substrates (melibiose, L-xylulose, phenylpropanoate, xanthosine/XMP) & FBA & \checkmark \\
C4c & GalNAc utilization is B2-advantaged; K-12 fails on GalNAc alone & FBA & \checkmark \\
C4d & Mucus glycans (GlcNAc, sialic acid, fucose) support \emph{E. coli} growth in FBA & FBA & \checkmark \\
C5 & B2 strains lack the \emph{frl} operon (frlA, B, C, D, R) & Genomic & \checkmark (LF82, UTI89, NRG857c) \\
C6 & LF82/UTI89/NRG857c are B2; K-12 MG1655 is A & Genomic & \checkmark (Clermont 2013 in-silico) \\
\bottomrule
\end{longtable}

\section{Method}
Three independent evidence layers, all from free public sources:
\begin{enumerate}
  \item \textbf{FBA on K-12 reference GEMs} (\texttt{iML1515}, \texttt{iJO1366}
    from BiGG). COBRApy 0.31.1, GLPK. Defined M9 minimal medium
    (NH$_4$, P$_i$, SO$_4$, K, Na, Mg, Ca, Fe$^{2+/3+}$, Cl, CO$_2$,
    H$_2$O, O$_2$, trace metals), all carbon exchanges closed, one test
    carbon opened at 10 mmol gDW$^{-1}$ h$^{-1}$. Ran every Table-1
    substrate and the five mucus-glycan monosaccharides (Fig.~3b).
  \item \textbf{Direct genomic verification on B2 reference genomes.}
    Downloaded LF82 (GCA\_000284495.1), UTI89 (GCA\_000013265.1),
    NRG857c (GCA\_000183345.1) and K-12 MG1655 (GCF\_000005845.2) via
    the free NCBI Datasets v2alpha REST API. Extracted K-12 \emph{frl}
    proteins (frlA/B/C/D/R) by accession, ran \texttt{tblastn} against
    each genome. Presence rule: pident $\geq$ 70\%, coverage $\geq$
    70\%, e-value $\leq 10^{-30}$. Also ran a 13-gene sanity panel of
    shared mucus-glycan / core catabolism genes.
  \item \textbf{In-silico Clermont phylogroup.} Encoded four Clermont
    (2013) quadruplex primers (chuA, yjaA, TspE4.C2, arpA). Required
    forward + reverse hits on the same contig, opposite strands, at
    expected amplicon distance. Mapped presence vector through Clermont
    2013 decision table.
\end{enumerate}

\section{Results vs.\ paper}

\subsection{Direct genomic verification of C5 (central B2 loss-of-function claim)}
\begin{tabular}{lccccccc}
\toprule
Strain & Phylogroup & frlA & frlB & frlC & frlD & frlR & Operon? \\
\midrule
LF82         & B2 & weak 29\% & weak 19\% & weak 23\% & weak 24\% & weak 23\% & \textbf{ABSENT} \\
UTI89        & B2 & weak 29\% & weak 19\% & weak 23\% & weak 25\% & weak 23\% & \textbf{ABSENT} \\
NRG857c      & B2 & weak 29\% & weak 19\% & weak 23\% & weak 24\% & weak 23\% & \textbf{ABSENT} \\
K-12 MG1655  & A  & 100\%/100\% & 100\%/100\% & 100\%/100\% & 100\%/100\% & 100\%/100\% & Present \\
\bottomrule
\end{tabular}

\medskip
Weak hits (19--29\% identity) in B2 strains are spurious cross-hits to
unrelated sugar transporters and kinases, far below the orthology
threshold. Clean, operon-wide loss.

\subsection{17-gene sanity panel}
Of 17 K-12 catabolism reference genes tested across the three B2 strains,
\textbf{16 conserved ($\geq$96\%) and only the 5-gene \emph{frl} operon
uniformly absent} in all three B2 strains: a textbook single operon-loss
signature in the B2 lineage.

\subsection{Independent phylogroup assignment via Clermont 2013}
\begin{tabular}{lcccccc}
\toprule
Strain & chuA & yjaA & TspE4.C2 & arpA & Signature & Call \\
\midrule
LF82        & + & $-$ & + & $-$ & \verb|+-+-| & B2 \\
UTI89       & + & $-$ & + & $-$ & \verb|+-+-| & B2 \\
NRG857c     & + & $-$ & + & $-$ & \verb|+-+-| & B2 \\
K-12 MG1655 & $-$ & + & $-$ & + & \verb|-+-+| & A \\
\bottomrule
\end{tabular}

\medskip
4/4 agreement with paper via a fully independent method.

\subsection{Table 1 substrates -- FBA on K-12 reference}
\begin{tabular}{lrrl}
\toprule
Substrate & iML1515 $\mu$ (h$^{-1}$) & iJO1366 $\mu$ (h$^{-1}$) & Match? \\
\midrule
Fructoselysine       & 0.893 & 1.005 & \checkmark \\
Psicoselysine        & 0.893 & 1.005 & \checkmark \\
Melibiose            & 1.770 & 1.967 & \checkmark \\
L-Xylulose           & 0.717 & 0.806 & \checkmark \\
Phenylpropanoate     & 1.131 & 1.259 & \checkmark \\
Xanthosine           & 1.023 & 1.214 & within variance \\
XMP                  & 1.023 & 1.214 & within variance \\
Cyanate (+ glucose)  & 0.886 & 0.993 & \checkmark (N-source) \\
\bottomrule
\end{tabular}

\medskip
6/8 qualitative match; 2 within reported within-phylogroup variance.

\subsection{Mucus-glycan FBA (Fig.~3b), K-12 reference}
GlcNAc 1.131, Neu5Ac 1.479, L-fucose 0.862, D-galactose 0.868,
D-glucuronate 0.705 (all $\mu$ in h$^{-1}$, iML1515). \textbf{GalNAc
alone: 0.000 (infeasible)} --- consistent with B2 TBP-aldolase advantage.

\subsection{GPR-level mechanism cross-check}
In iML1515, GPRs of the fructoselysine/psicoselysine pathway map onto
\emph{exactly} the genes we show are missing from B2:
\texttt{FRULYSt2pp} $\to$ \emph{frlA};
\texttt{FRULYSDG} $\to$ \emph{frlB};
\texttt{FRULYSK} $\to$ \emph{frlD};
\texttt{FRULYSE} $\to$ \emph{frlC};
\texttt{PSCLYSt2pp} $\to$ \emph{frlA}.
Chain closed end-to-end.

\section{Coverage / Agreement}
\textbf{Coverage: 8/10.} C1, C2, C4a--d, GPR mechanism, direct
genomic verification of C5, and C6 (Clermont). Outstanding: per-strain
GEMs for 53 IBD isolates, 649-substrate FBA panel, Fig.\ 1a/3a heatmaps.
\textbf{Agreement: 10/10.} Every tested Table-1 substrate behaved as the
paper predicts; entire \emph{frl} operon absent in every B2 reference;
all 4 Clermont calls match; all genome sizes match. No fluxes fabricated
(\texttt{cobra.Model.optimize()} outputs) and no BLAST hits invented.

\section*{GENUINE CRITIQUE}
\addcontentsline{toc}{section}{Genuine critique}

This section states honestly what did \emph{not} replicate and what
weaknesses remain --- both in the original paper and in this replication.

\subsection*{What did NOT replicate in this pass}
\begin{itemize}
  \item \textbf{The 110-strain pan-genome is not reproduced.} Only 4
    genomes (3 B2 + 1 A) were touched at the sequence level. Fig.\ 1a
    (110$\times$reactions matrix) and Fig.\ 3a (110$\times$649
    substrate-growth heatmap) remain unverified.
  \item \textbf{The 53 IBD-patient isolates were not individually
    modeled.} The paper's over-representation statistic for B2 in IBD
    (its epidemiological punchline) rests on those 53 strains; we did
    not remap the strain-to-BV-BRC-accession list or rebuild their GEMs.
  \item \textbf{The 649-substrate FBA panel was not run.} Only 8
    Table-1 substrates and 7 mucus glycans were tested; the paper's
    Fig.\ 3a scores 649 carbon sources per strain.
  \item \textbf{No per-strain GEM was reconstructed.} FBA used the K-12
    reference iML1515/iJO1366, which the paper puts in phylogroup A.
    Growth-rate predictions for the 53 B2 isolates are therefore
    inferred from the K-12 model + paper framework, not run directly.
  \item \textbf{Two Table-1 substrates (xanthosine, XMP) partially
    disagree.} K-12 grows on both at $\mu \approx 1.0$--$1.2$ h$^{-1}$
    while the paper has only 38--46\% of non-B2 strains grow. This is
    within the paper's reported within-phylogroup variance but a clean
    REPLICATED tag would require matching the growing \emph{fraction},
    not just K-12's binary answer.
\end{itemize}

\subsection*{Weaknesses in the original paper}
\begin{itemize}
  \item \textbf{Sample size is modest for a phylogroup-level claim.} 53
    IBD isolates spread across 7 phylogroups, so per-phylogroup n in
    the IBD cohort is small (single or low double digits). The
    ``B2 over-represented in IBD'' inference is population-level rather
    than mechanistic within-host.
  \item \textbf{Correlation vs.\ causation.} That B2 lacks \emph{frl}
    and B2 is over-represented in IBD does not by itself show that
    \emph{frl} loss \emph{causes} the association. The paper's own
    strongest claim is about metabolic capability, not causal
    contribution to IBD pathogenesis.
  \item \textbf{Cross-sectional design.} No longitudinal within-patient
    tracking of B2 dominance during flare vs.\ remission; the metabolic
    argument is built on cross-cohort comparison.
  \item \textbf{FBA growth $\neq$ in-vivo fitness.} Table-1 growth-rate
    predictions in defined M9 do not directly measure competitive
    fitness in the mucus layer, where oxygen, redox, cross-feeding, and
    host defenses all matter.
  \item \textbf{Reference-strain bias in the mucus-glycan story.} The
    TBP-aldolase / GalNAc advantage argument leans on gene-content
    comparisons; per-strain wet-lab growth on GalNAc is not shown in
    the paper for all 53 IBD isolates.
\end{itemize}

\subsection*{Weaknesses in this replication}
\begin{itemize}
  \item Only 3 B2 genomes tested (LF82, UTI89, NRG857c); cannot speak
    to within-B2 variance.
  \item Presence rule (pident $\geq$70, cov $\geq$70, e $\leq10^{-30}$)
    is a reasonable comparative-genomics default but not the paper's
    exact CD-HIT-at-80\% clustering; a strict method-match rerun could
    move edge cases.
  \item Clermont in-silico PCR allows $\leq$15\% primer mismatch and is
    generous on amplicon-size windows; would disagree with the paper's
    pan-genome-based phylogrouping on edge cases (none in our 4-strain
    panel, but not stress-tested).
  \item BV-BRC strain-name mapping for the 53 IBD isolates was not
    completed.
\end{itemize}

\subsection*{What would move this to full REPLICATED}
\begin{enumerate}
  \item Download all 110 genome FASTAs from Additional file 1 /
    Table S1.
  \item Run CD-HIT at 80\% identity across 110 proteomes; rebuild the
    pan-genome.
  \item Reconstruct per-strain GEMs via CarveMe or KBase (${\sim}$1--2
    weeks single-workstation compute).
  \item Loop the 649-substrate panel through
    \texttt{cobra.Model.optimize()} per strain; rebuild Figs.~1a/3a.
  \item Cross-check IBD-cohort epidemiological statistics against the
    53-strain metadata as reported.
\end{enumerate}
None of these steps require commercial software, GPUs, or restricted
data.

\section{Resources}
Europe PMC, BV-BRC public API, NCBI Datasets v2alpha REST (free, no
auth), BiGG (iML1515, iJO1366), COBRApy 0.31.1 + GLPK, BLAST+
(\texttt{makeblastdb}, \texttt{tblastn}, \texttt{blastn}), Biopython
1.87. Total compute: ${\sim}$3 min laptop wall-clock.

\section{Reproducibility}
See \texttt{work/} for scripts and outputs:
\texttt{genome\_stats.py, frl\_blast.py, metabolic\_survey.py,
clermont.py, fba\_table1.py, fba\_mucus.py} and corresponding
\texttt{.json} outputs; BLAST databases under \texttt{blast/};
downloaded assemblies under \texttt{genomes/}.

End-to-end reproduction:
\begin{verbatim}
python3 -m venv .venv && source .venv/bin/activate
pip install cobra biopython
for acc in GCA_000284495.1 GCA_000013265.1 GCA_000183345.1 GCF_000005845.2; do
  curl -sS -o "${acc}.zip" \
   "https://api.ncbi.nlm.nih.gov/datasets/v2alpha/genome/accession/${acc}/download\
?include_annotation_type=PROT_FASTA&include_annotation_type=GENOME_FASTA"
  unzip -q "${acc}.zip" -d "${acc}"
done
python3 genome_stats.py
python3 frl_blast.py
python3 metabolic_survey.py
python3 clermont.py
python3 fba_table1.py
python3 fba_mucus.py
\end{verbatim}

\end{document}
